\subsection{The Eucharist: Sacrifice, Presence, Communion, and Pledge}

\begin{studybox}{Governing thesis}
The Most Holy Eucharist is the sacrament in which the one sacrifice of Christ is made sacramentally present and offered under the species of bread and wine; by consecration their whole substances become his Body and Blood while the sensible species remain. Christ, whole and entire, is the \emph{res et sacramentum}; union in charity with Christ and his ecclesial Body is the further \emph{res}. Because it contains the author of grace himself, the Eucharist completes initiation, gathers every sacrament toward the Pasch, and pledges bodily resurrection and the wedding supper of the Lamb.
\end{studybox}

\subsubsection{Creation offered back: bread, wine, altar, and banquet}

Bread and wine begin as created gifts and humanly cultivated fruits. Grain is gathered, broken, mixed, and baked; grapes are harvested and pressed. Their production joins earth, weather, time, intelligence, labor, and community. The Eucharist therefore does not descend as an escape from creaturely mediation. The incarnate Son receives creation through thanksgiving and returns it to the Father as bearer of his own self-gift. Nature is not annihilated but taken beyond every native potency by divine conversion.

Abel's acceptable first fruits and Noah's altar place thanksgiving and sacrifice near humanity's beginnings (Gen. 4:3--5; 8:20--22). Melchizedek, priest of God Most High and king of Salem, brings bread and wine and blesses Abram (Gen. 14:18--20). Hebrews' principal use of Melchizedek concerns Christ's eternal priesthood, not a simplistic identity between the Genesis offering and the Mass; patristic tradition nevertheless rightly sees bread, wine, priesthood, blessing, and peace converging providentially toward the Eucharistic King.

Passover gives the controlling covenant pattern. An unblemished lamb is slain, its blood marks deliverance from death, its flesh is eaten, and the memorial makes the saving act liturgically present to later generations (Exod. 12). Israel then passes through the sea and receives manna in the wilderness. At Sinai sacrifice, proclaimed word, blood, assent, and a meal before God ratify covenant communion (Exod. 16; 24:3--11). The Eucharist fulfills these without becoming one more annual rescue: Christ is Lamb, liberating blood, true food, covenant mediator, and leader of the final Exodus.

The bread of the Presence remains before the Lord; cereal and wine offerings join Israel's worship; the peace offering culminates in a sacred meal (Lev. 2; 7; 24:5--9). Wisdom sets a table and summons the simple to eat her bread and drink her wine; Elijah is strengthened by bread for a journey beyond his strength (1~Kgs. 19:4--8; Prov. 9:1--6). Malachi foresees a pure offering among the nations from rising to setting of the sun (Mal. 1:11). The Fathers read this catholic sacrifice in the Church's Eucharist. The prophecy's moral demand remains: a pure offering cannot be invoked to excuse contempt for the poor, faction, or idolatry.

\begin{fourstable}{Old-covenant sign}{Christological fulfillment}{Eucharistic realization}{Limit on the type}
Melchizedek's bread, wine, priesthood, and blessing & Christ is royal priest by indestructible life, not Levitical descent. & The Church offers bread and wine that Christ converts into his sacrificial presence. & Hebrews emphasizes person and priesthood; Genesis alone is not a complete Eucharistic rite.\\
Passover lamb and blood & Christ the unblemished Lamb passes through death and liberates from sin. & His Body is eaten and covenant Blood received in memorial sacrifice. & The Mass neither kills Christ again nor merely reenacts an absent event.\\
Manna and water in the wilderness & The Father gives the true Bread from heaven and living water in the Spirit. & Pilgrims receive food containing the risen Lord and a pledge of immortality. & Israel's recipients still died; the Eucharist bears fruit only in living communion.\\
Sinai sacrifice and banquet & The Son mediates the new covenant by obedient self-offering. & Word, consent, sacrifice, Blood, and ecclesial meal form one liturgical act. & Communion cannot be severed from covenant obedience and reconciliation.\\
\end{fourstable}

\subsubsection{Signs multiplied and discourse in Capernaum}

Jesus' table ministry already reveals messianic abundance, mercy toward sinners, and the danger of receiving fellowship without conversion. At Cana he gives superabundant wine as the first sign of nuptial glory; in the wilderness he multiplies loaves, gives thanks, breaks, and distributes through disciples (John 2:1--11; 6:1--15; Mark 6:30--44). The signs are not proofs by resemblance alone. They disclose the identity and authority of the giver whose hour will convert wedding joy and wilderness food into the gift of his Paschal flesh.

John 6 moves from sign to faith and from faith to eating. Jesus is the true bread because he comes from the Father and gives life; unbelief cannot be cured by material appetite. Yet the discourse intensifies rather than retreats: the bread he will give is his flesh for the life of the world; his flesh is true food and his blood true drink; the one who eats and drinks abides in him and will be raised on the last day (John 6:32--58). Spirit and life do not cancel flesh. In John, the Spirit makes Christ's life-giving flesh available; a merely carnal understanding would imagine ordinary torn flesh apart from Paschal gift and sacramental mode.

Many disciples withdraw, while Peter remains because Jesus has the words of eternal life (John 6:60--69). Eucharistic faith is therefore neither credulity about appearances nor mastery of a mechanism. It trusts the incarnate Word precisely where senses report bread and wine and divine speech gives a deeper judgment. The promise of resurrection shows the final cause: Eucharistic eating joins the mortal body to the risen flesh of the Son.

\subsubsection{The Supper interprets the Cross}

The Synoptic Gospels and Paul transmit the institution in four closely related accounts (Matt. 26:26--29; Mark 14:22--25; Luke 22:14--20; 1~Cor. 11:23--26). Jesus takes bread, blesses or gives thanks, breaks, gives, and says it is his Body given ``for you.'' Over the cup he identifies his Blood of the covenant, poured out for many for remission. He commands, ``Do this in memory of me,'' and points toward the kingdom's future banquet.

These words interpret the next day's death before it occurs. The distinct consecration of Body and Blood signifies sacrificial separation; the living Christ is not physically divided under the species, but his self-offering unto death is sacramentally represented. The Supper anticipates Calvary sacramentally; every Mass thereafter makes the same completed sacrifice present under signs. Historical sequence differs, sacrificial identity does not.

Biblical memorial, \emph{anamnesis}, is more than mental recollection. Passover liturgy places later Israel within God's saving deed; Christ commands an ecclesial action in which his once-for-all offering becomes operative now. The Church does not drag the glorified Lord back into suffering. The risen Lord makes his unrepeatable act present in the mode of sacrament and applies its fruits to living and dead.

The covenant words evoke Exodus 24 and Jeremiah's promised new covenant. Christ is simultaneously priest, victim, altar in the profound Christological sense, covenant mediator, and food. ``For you'' makes sacrifice gift before it becomes the disciples' offering; ``Do this'' establishes apostolic service; eating and drinking make objective presence tend toward communion. No one aspect is sufficient alone.

\subsubsection{Paul: participation, one Body, and judgment}

Paul asks whether the cup of blessing is a participation---\emph{koinonia}---in Christ's Blood and the broken bread a participation in his Body. Because there is one Bread, the many are one body (1~Cor. 10:16--17). He compares Israel's sacrificial meals and pagan altars to show that eating establishes a real cultic communion. The argument would collapse if Eucharistic participation were a mental recollection devoid of objective referent.

In Corinth social contempt contradicts the sacrament. Wealthier Christians eat ahead while others go hungry; Paul responds by handing on the institution and warning that unworthy eating and drinking profane the Lord's Body and Blood, bring judgment, and require discernment (1~Cor. 11:17--34). ``Unworthily'' describes the manner and disposition, not a demand that communicants first become intrinsically worthy of Christ. Grace is gift, but a gift of unity is received against its meaning when grave sin, faction, or contempt is knowingly embraced.

The Eucharist thus judges the Church it makes. Sacramental validity does not automatically create the final unity where recipients place an obstacle. Yet human failure cannot reduce Christ's Body to ordinary bread: Paul's warning derives its seriousness from objective presence. The same realism grounds adoration and ethics. One cannot honor the Head under sacramental veils and despise his members in the assembly, the poor, the unborn, the stranger, or the suffering.

Luke's Emmaus narrative gathers word, Paschal interpretation, hospitality, blessing, breaking, recognition, disappearance, and mission (Luke 24:13--35). Acts describes Christians persevering in apostolic teaching, fellowship, breaking of bread, and prayers (Acts 2:42). These texts do not prove that every biblical ``breaking of bread'' is technically a Mass; they reveal the Eucharistic shape of apostolic life: Scripture opened by the risen Lord, table recognition, communion, and return to the assembled Church.

\subsubsection{Patristic realism from Antioch to Hippo}

Ignatius of Antioch says Docetists abstain from Eucharist and prayer because they do not confess the Eucharist to be the flesh of Jesus Christ that suffered for sins and was raised by the Father (\emph{Smyrnaeans} 7). Eucharistic realism is here Christological: denial of real flesh in the sacrament accompanies denial of real Incarnation and Passion. Ignatius also binds a secure Eucharist to bishop or authorized celebrant, making presence inseparable from visible ecclesial communion (\emph{Smyrnaeans} 8).

Justin describes Sunday worship: readings, exhortation, common prayers, kiss, bread, wine and water, presidential thanksgiving, the people's Amen, distribution by deacons, and aid for the needy (\emph{First Apology} 65--67). The food is not received as common bread or drink, but as the flesh and blood of the incarnate Jesus through the prayer derived from him. His description witnesses word, anaphora, real presence, ordained presidency, diaconal distribution, Sunday resurrection, and charity without supplying a later scholastic analysis.

Irenaeus argues against Gnostics from the Eucharist. Bread from earth, receiving God's invocation, is no longer common bread but Eucharist, consisting of earthly and heavenly realities; bodies nourished by Christ's Body and Blood are destined for resurrection (\emph{Against Heresies} IV.18.4--5; V.2.2--3). The Creator's gifts can become instruments of immortality because the Word truly assumed and saves flesh. Sacramental change therefore refutes contempt for matter rather than confirming it.

Cyril of Jerusalem instructs neophytes not to judge by taste alone: Christ's declaration makes bread his Body and wine his Blood (\emph{Mystagogical Catechesis} IV.1--6). Ambrose asks whether Christ's word, which created what did not exist, can change existing natures; he places consecration within the Word's creative authority (\emph{De mysteriis} 9.50--54). Neither Father uses Aristotle's later technical vocabulary, but both confess an objective conversion grounded in divine speech.

John Chrysostom speaks with awe of the same Christ who furnished the Supper's table now bringing about the mystery, and draws Eucharistic communion toward care for Christ in the poor (\emph{Homilies on Matthew} 82; \emph{Homilies on First Corinthians} 24, 27). Augustine can call the Eucharist a sign and also insist upon its reality. His preaching to the baptized moves from the many grains made one loaf to the summons to become Christ's ecclesial Body by receiving his sacramental Body (Sermons 227, 272; \emph{Tractates on John} 26). Patristic ``symbol'' means a mystery that communicates its reality, not an empty pointer.

Ephrem's Syriac hymns approach the mystery through fire and Spirit hidden in bread and wine, training reverent wonder where discursive definition reaches its limit. Eastern emphasis upon epiclesis and transformation and Western emphasis upon Christ's determining words do not name competing agents. The one Trinitarian action is addressed and articulated differently within integral anaphoras.

\begin{fourstable}{Witness}{Controlling claim}{Doctrinal fruit}{Misuse excluded}
Ignatius, \emph{Smyrnaeans} 7--8 & Eucharist is Christ's suffering and risen flesh within episcopal communion. & Presence, Incarnation, sacrifice, and visible Church cohere. & Realism cannot be detached from ecclesial unity or used as mere anti-symbol slogan.\\
Justin, \emph{Apology} I.65--67 & Consecratory thanksgiving distinguishes Eucharist from common food. & Word, anaphora, ordained ministry, reception, Sunday, and almsgiving form one act. & His outline is not a verbatim universal missal.\\
Irenaeus, \emph{Against Heresies} IV.18; V.2 & Created gifts become Eucharist and nourish bodies destined to rise. & Eucharistic doctrine defends creation, Incarnation, and bodily resurrection. & ``Earthly and heavenly'' does not leave unchanged bread alongside Christ.\\
Cyril, Ambrose, Augustine & Divine declaration converts; senses yield to faith; sacramental Body makes ecclesial Body. & Objective presence tends to charity and unity. & ``Sign'' and ``figure'' in patristic usage do not mean absence of the signified reality.\\
\end{fourstable}

\subsubsection{One sacrifice sacramentally present}

Hebrews insists that Christ offered himself once for all and, unlike repeated Levitical victims, has perfected his people by one sacrifice (Heb. 7:27; 9:12, 24--28; 10:10--14). Catholic doctrine does not place the Mass on the repeated side of Hebrews' contrast. The Eucharistic sacrifice is numerically one in principal priest and victim with Calvary; the manner of offering differs. On the Cross Christ offered bloodily and acquired redemption; in the Mass the risen Christ offers through sacramental signs and applies the fruits of that completed sacrifice.

The ordained priest acts \emph{in persona Christi Capitis} when pronouncing Christ's words and offering the sacrifice, but Christ remains principal priest. The whole Church also offers in a differentiated way: the baptized unite themselves, praise, labor, suffering, and intercession to the Head's offering. Ministerial and baptismal priesthood are ordered to one another without becoming the same power. Only bishop or presbyter consecrates; every faithful communicant is summoned to become an offered life (Rom. 12:1; 1~Pet. 2:5).

As sacrifice, the Mass is first a Godward act of religion: through Christ the Church renders the Father adoration and glory, praise and thanksgiving. Because the sacramentally present victim is the same obedient Christ whose Blood remits sins, this one oblation is also propitiatory and impetratory, and may be offered for the living and the dead. These applications are subordinate to the sacrifice's Godward finality without being optional or unreal (ST II--II, q.~85, aa.~1--3; ST III, q.~83, a.~1; Trent, sess.~XXII, chs.~1--2, can.~3; CCC 1359--1361, 1367--1371). Human payment does not purchase grace, the Mass supplies no defect in the Cross, and its fruits are received according to divine wisdom and creaturely disposition.

Sacrifice and banquet cannot be separated. Consecration constitutes sacramental sacrifice even if no layperson communicates, while the celebrating priest must communicate the victim whose offering he completes ritually. The faithful's Communion is a fuller participation and the sacrament's intrinsic orientation. Conversely, a meal account without altar, covenant Blood, and self-offering cannot explain the institution or Paul's sacrificial comparisons.

\subsubsection{Transubstantiation: conversion beneath appearances}

The Church confesses that after consecration Christ is truly, really, and substantially present and that the whole substances of bread and wine are converted into the substances of his Body and Blood; the appearances or species remain. Trent judges ``transubstantiation'' a fitting and proper name for this unique conversion (sess.~XIII, ch.~4). The term protects mystery from two reductions: bread does not merely acquire a sacred use while remaining bread, and Christ does not become locally enclosed as another visible object beside it.

Aristotle's distinction between substance and accidents gives disciplined language. Substance answers what a thing is; accidents include sensible qualities, dimensions, taste, weight, and chemical behavior. In ordinary substantial change, new substance bears new accidents. Here divine power converts the underlying substances while sustaining the accidents without their former subjects (ST III, qq.~75, 77). The doctrine does not depend upon an obsolete guess about microscopic physics. Scientific analysis measures the sacramental species precisely in the order where appearances remain; it cannot assay substantial presence as though Christ were another chemical constituent.

The conversion is total and terminates in Christ's pre-existent glorified Body and Blood; it is not annihilation followed by replacement, local motion from heaven, or generation of a new Christ (ST III, q.~75, aa.~3--4, 7). His presence is sacramental rather than circumscribed. He is whole under each species and whole under every part after division, while the species endure. Breaking divides appearances, not Christ's body.

By the force of the words, Body is present under the species of bread and Blood under wine; by natural concomitance, because the risen Christ can die no more, Blood, soul, and divinity accompany the Body, and Body, soul, and divinity accompany the Blood. The separate species represent sacrificial death without physically separating the living Lord. Communion under either species therefore gives whole Christ and the full grace of the sacrament, though reception under both more fully expresses the banquet's sign when law provides.

Presence endures as long as the Eucharistic species remain. Reservation, adoration, genuflection, procession, care for fragments, and Communion of the sick follow from the sacrament's identity, not from later replacement of an original meal. When the species are corrupted so that bread or wine no longer remain perceptibly, sacramental presence ceases; Christ does not undergo digestion, injury, or corruption.

\begin{studybox}{Metaphysical guardrail}
Transubstantiation does not ask the senses to report what lies outside their proportion. Sense judges species truly; faith, resting on Christ's word, judges substance. Nor does ``substance'' mean a hidden physical particle. The conversion is ontological and sacramental, caused by divine power, with the appearances sustained so that Christ can be received as food without local displacement or bodily mutilation.
\end{studybox}

\subsubsection{Matter, form, and the integrity of the anaphora}

Valid matter is bread made from wheat and natural wine of the grape. In the Latin Church the priest must use unleavened bread; most Eastern Churches use leavened bread, while Armenians also use unleavened bread. Leaven concerns lawful ritual patrimony, not two different Eucharists. The wine must not be corrupted; a small quantity of water is mixed by ecclesial law and rich symbolism, but water is not part of the substantial matter required for validity (CIC 924--926; CCEO 706--707).

The 1962 Roman Missal places the determining words within the one Roman Canon. Over the host the form is \latin{Hoc est enim Corpus meum}. Over the chalice it is \latin{Hic est enim Calix Sanguinis mei, novi et aeterni testamenti: mysterium fidei: qui pro vobis et pro multis effundetur in remissionem peccatorum}. The words must be spoken by a priest intending to consecrate this valid matter within the sacrificial act. The punctuation is editorial; sacramental identity lies in the substantial sign, not typography.

The current Roman Missal retains the institution narrative but places \latin{mysterium fidei} as the invitation to the memorial acclamation rather than within the chalice words. An exposition for the 1962 books should quote its own text accurately and cross-note, not conflate, the present form. Neither form is a private script for detached recitation; the liturgical Eucharistic Prayer belongs to validity's ecclesial context and to lawful celebration's integrity.

Eastern anaphoras characteristically display the invocation of the Holy Spirit, institution narrative, memorial, and oblation as an integral consecratory movement. Catholic theology need not choose between Christ's word and the Spirit's operation: the Son's command is effective in the Spirit through the prayer of his Body. The Holy See's judgment concerning the ancient Anaphora of Addai and Mari is an essential guardrail. Although it lacks a coherent ad-litteram institution narrative, the words of institution are present in a dispersed euchological way within a valid apostolic anaphora. One must not turn a Latin scholastic extraction into a test that contradicts the Church's authoritative judgment about an Eastern liturgy.

\begin{massbox}{Causal constitution of the Eucharist}
\textbf{Matter:} wheaten bread and natural grape wine, offered according to an approved rite. \textbf{Form:} Christ's consecratory word within the Church's integral Eucharistic Prayer / anaphora, articulated according to the rite. \textbf{Principal cause and priest:} Christ, whose glorified humanity makes the one Pasch present. \textbf{Instrumental minister:} a validly ordained bishop or presbyter alone. \textbf{Intermediate reality:} whole Christ under the enduring species. \textbf{Sacrificial end:} Godward adoration and thanksgiving, with subordinate propitiation and impetration. \textbf{Communion's proper end:} union with Christ and his Body in charity. \textbf{Ultimate end:} resurrection and the heavenly banquet in unveiled communion.
\end{massbox}

\subsubsection{Minister, communicant, and disposition}

A validly ordained bishop or presbyter alone confects the Eucharist (CIC 900; CCEO 699). There is no extraordinary consecrator. A deacon is an ordinary minister of distributing Holy Communion in Latin law; an instituted acolyte or deputed faithful person can be an extraordinary minister of Communion. Those offices concern reverent distribution of a sacrament already present and cannot be described as ``performing the Eucharist.'' In Eastern discipline the terminology and order of distribution follow the competent law.

The Eucharist is constituted by consecration independently of a later recipient, but sacramental Communion can be received only by the baptized who are not prohibited by law. A person conscious of grave sin must ordinarily receive absolution before Communion; where grave reason exists and confession is unavailable, perfect contrition including resolve to confess as soon as possible is required (CIC 916). The Eucharist forgives venial sins and preserves from future mortal sin by strengthening charity; it is not ordered to forgive unrepented mortal sin as Baptism or Penance is.

Latin children receive when capable of distinguishing the Eucharist from ordinary food and receiving with faith and devotion, with special provision for danger of death; Eastern infants receive according to their liturgical tradition (CIC 913--914; CCEO 710). The difference is disciplinary and mystagogical, not a disagreement about whether infants are baptized members capable of grace. Adults should observe the Eucharistic fast and the competent Church's frequency and reception law. Historical 1962 fasting discipline and current canon law must be labeled by era rather than blended.

Spiritual Communion is an act of desire, faith, and charity ordered toward sacramental reception. It can bear abundant grace where reception is impossible or presently unfitting, but it does not consecrate, place the sacramental species in the recipient, or make ecclesial and moral preparation optional. Likewise, viewing a liturgy through media can assist prayer but is not bodily sacramental attendance or Communion.

\subsubsection{Sacramental operation and the hierarchy of ends}

In Aquinas's Eucharistic application, the consecrated species are the \emph{sacramentum tantum}; Christ's true Body and Blood are the \emph{res et sacramentum}, effected under the sign and signifying ecclesial unity; the further \emph{res tantum} is grace and the unity of the Mystical Body (ST III, q.~73, a.~6). The intermediate reality here is not character. Christ is substantially present whether or not a particular communicant receives fruitfully.

The Eucharist must be considered under two inseparable operations. In the Mass as sacrifice, Christ and the Church act Godward by offering the one Paschal victim; in sacramental Communion, Christ nourishes the recipient into union with himself and therefore with his Mystical Body. The threefold sacramental structure identifies sign, contained reality, and final sacramental grace; the following hierarchy identifies the governing and subordinate ends of sacrifice and reception without turning their fruits into a flat catalogue.

\begin{fourstable}{Ordered level}{Eucharistic fulfillment}{Teleological status}{Controlling loci}
\emph{Sacramentum tantum} & The enduring species of bread and wine are the outward sacramental sign under which consecration, offering, adoration, and reception occur. & The species truly signify food, sacrifice, and the unity of many in one; they are not bread and wine remaining in substance after consecration. & ST III, q.~73, a.~6; q.~75, a.~4; Trent, sess.~XIII, chs.~3--4; CCC 1373--1377.\\
\emph{Res et sacramentum} & Whole Christ---Body, Blood, soul, and divinity---is substantially present under each species and every part while the species endure. & Christ is the reality contained and the sign of ecclesial unity, objectively present before fruitful reception; unlike Baptism or Confirmation, this intermediate reality is not character. & ST III, q.~73, a.~6; q.~76, aa.~1--3; Trent, sess.~XIII, chs.~1, 3; CCC 1374, 1377.\\
\emph{Res tantum} & Grace and union with Christ in the unity of his Mystical Body: the worthy communicant lives from the Head and is incorporated more deeply into his members. & This final sacramental reality is caused through reception; it must be distinguished from Christ's objective presence and can be impeded by contrary disposition. & ST III, q.~73, a.~6; q.~79, aa.~1, 8; Trent, sess.~XIII, chs.~2, 7; CCC 1391, 1396.\\
Sacrifice: primary Godward ends & The Church, united to Christ's self-offering, renders latria, praise and glory, and thanksgiving to the Father through the sacramental memorial of the one Pasch. & Godward worship is constitutive and governing: sacrifice is offered to God, not primarily to produce a benefit in the offerer or recipient. & ST II--II, q.~85, aa.~1--3; ST III, q.~83, a.~1; Trent, sess.~XXII, ch.~1; CCC 1359--1361, 1367--1369.\\
Sacrifice: subordinate applications & The same oblation is propitiatory and impetratory for the living and the dead, sacramentally applying the satisfaction and intercession of the one victim. & Propitiation and petition are objective ends of the offering, not optional intentions; their determinate fruits do not operate mechanically but vary with disposition and providence. & ST III, q.~79, aa.~5, 7; Trent, sess.~XXII, ch.~2, can.~3; CCC 1368--1371.\\
Communion: primary proper / proximate end & Worthy sacramental eating unites the communicant to Christ in grace and charity and, inseparably, more deeply to the one Mystical Body. & This union governs Communion's operation. Ecclesial unity is neither a merely social aftereffect nor separable from personal adhesion to the Head. & ST III, q.~73, a.~6; q.~79, a.~1; Trent, sess.~XIII, ch.~2; CCC 1391, 1396.\\
Communion: subordinate proper effects & Grace and charity increase; venial sins are remitted; disordered attachments are healed; future sin is resisted; spiritual delight is given; glory and bodily resurrection are pledged. & These are intrinsic consequences ordered beneath union. Remission is not the principal fruit, and no sensible intensity or particular outcome displaces union with Christ and his Body. & ST III, q.~79, aa.~1--2, 4--6, 8; Trent, sess.~XIII, chs.~2, 7, can.~5; CCC 1392--1395, 1402--1405.\\
Common ultimate end & Perfect charity in the risen Body, the wedding supper of the Lamb, and face-to-face possession of God when sacramental species yield to vision. & Eucharistic sacrifice and Communion converge in beatitude: God is perfectly glorified in Christ and the redeemed share the Son's filial communion. & ST III, q.~65, a.~3; q.~79, a.~2; CCC 1130, 1402--1405.\\
\end{fourstable}

The primary proper effect in a worthy communicant is union with Christ and, through the Head, his Body. Bodily food is assimilated into the eater; this food, by a higher causality, assimilates the eater to Christ. Increase of grace and charity, venial remission, healing of attachments, preservation from future sin, spiritual delight, and the pledge of resurrection are real subordinate or consequent effects (ST III, q.~79, aa.~1--2, 4--6, 8; CCC 1391--1405). ``Subordinate'' ranks these fruits beneath union; it does not deny their efficacy, while their fruitful realization must still be distinguished from the sacrament's objective constitution.

The sacrament's ecclesial effect is demanding. Augustine's many grains and one loaf interpret Paul's one Bread and one Body. Communion is not a private possession of ``my Jesus'' alongside indifference to doctrine, hierarchy, neighbor, or poor. Neither is ecclesial unity manufactured sociologically by eating together. Christ's prior Body makes members one and calls their free acts into conformity with the received reality.

Frequent and even daily Communion can be fitting where right disposition and law are observed; repetition does not imply the first Eucharist was insufficient. Food sustains a pilgrim life through repeated acts, and every Mass sacramentally applies the inexhaustible Pasch. Unlike Baptism, Confirmation, and Orders, the Eucharist imprints no character because it gives the end toward which sacramental deputation tends: union with Christ himself (ST III, q.~65, a.~3).

\subsubsection{Center of the sevenfold organism}

Baptism makes a member capable of the table; Confirmation seals the member for the worship and witness flowing from it. Penance restores communion when grave sin contradicts baptismal life. Anointing joins sick flesh to the sacrifice and prepares for Viaticum. Orders exists chiefly to make Christ's Eucharistic action available and to shepherd the Body it forms. Matrimony signifies the Bridegroom's covenant and is nourished by the sacrificial self-gift it represents. Thus the Eucharist is not simply fourth in a list; it is initiation's summit and the center toward which the other sacraments tend (ST III, q.~65, a.~3; CCC 1324).

This ordering also prevents Eucharistic isolation. A consecration detached from baptismal faith, apostolic ministry, reconciliation, service, and nuptial ecclesiology would contradict the economy that gives it meaning. Conversely, every Christian doctrine tends toward doxology: Trinity, Incarnation, Cross, resurrection, ascension, Spirit, Church, grace, body, creation, judgment, and glory meet at the altar without being exhausted by it.

The Eucharist is medicine of immortality, yet communicants still die. Sacramental grace works in pilgrim mode, restoring and strengthening a body still awaiting transformation. Viaticum makes the final direction explicit: Christ accompanies the dying baptized through the last passage as food of resurrection. The last earthly Communion is not a charm against physical death but pledge that death cannot sever the member who abides in the risen Head.

\subsubsection{Until he comes}

Paul says the Church proclaims the Lord's death ``until he comes'' (1~Cor. 11:26). Eucharistic presence is real and therefore intensifies, rather than abolishes, eschatological hunger. Christ is possessed under sacramental veils, not yet seen face to face. The liturgy joins angels and saints while the pilgrim still prays for deliverance, peace, resurrection, and admission into their company.

Revelation's wedding supper of the Lamb and the holy city adorned as a bride disclose the Eucharist's final cause (Rev. 19:6--9; 21:1--4, 22--27). The altar prepares a people to become the Bride: washed, sealed, fed, reconciled, healed, ordered, and made fruitful. In glory there is no temple as a separate sacred enclosure, for God and the Lamb are its temple; no sacramental species mediate presence, for the servants see his face. The charity formed through Eucharistic Communion remains when sacramental manner yields to vision.

\begin{studybox}{Guardrails for Eucharistic speech}
The Mass does not kill Christ again; memorial is not mental recall; ``symbol'' does not mean absence; transubstantiation is not a chemical claim; Christ is not physically torn when the host is broken; Communion under one species is not half of Christ; extraordinary ministers of Communion cannot consecrate; infant Eastern Communion is not a defective discipline; and the integral Eastern anaphora cannot be judged by silently treating one Roman textual arrangement as universal.
\end{studybox}

\begin{massbox}{The Eucharist in the sacramental whole}
\textbf{From:} creation's fruits, Melchizedek, Passover, manna, covenant sacrifice, and promised banquet. \textbf{Through:} the incarnate Word's once-for-all Pasch made present by priestly consecration under bread and wine. \textbf{As sacrifice:} adoration, praise, thanksgiving, propitiation, and petition. \textbf{As Communion:} union with Christ in charity and the unity of his Body. \textbf{Unto:} bodily resurrection, the wedding supper, and face-to-face possession when sacramental veils pass away.
\end{massbox}
